\documentclass[12pt]{article}
\usepackage{amsmath,amssymb,amsthm,hyperref,booktabs,xcolor,enumitem,listings}
\usepackage[margin=1in]{geometry}

\newtheorem{theorem}{Theorem}
\newtheorem{lemma}[theorem]{Lemma}
\newtheorem{definition}[theorem]{Definition}
\newtheorem{remark}[theorem]{Remark}
\newtheorem{corollary}[theorem]{Corollary}

\definecolor{proved}{RGB}{40,120,40}
\definecolor{open}{RGB}{160,50,50}
\definecolor{codeblue}{RGB}{30,60,140}

\lstset{
  basicstyle=\ttfamily\small,
  keywordstyle=\color{codeblue},
  breaklines=true,
  frame=single
}

\title{\textbf{T\_EML\_WFP\_ANALYTIC: Lock Document}\\[0.4em]
\large Well-Formed EML Trees are Real-Analytic on $(0,\infty)$\\
\normalsize Theorem 25 of the EML Verification Sprint}
\author{Arturo R.\ Almaguer}
\date{2026-04-21}

\begin{document}
\maketitle

\begin{abstract}
We state, prove, and lock \textbf{T\_EML\_WFP\_ANALYTIC}: every EML tree satisfying the
\emph{WellFormedPos} condition is real-analytic on $(0,\infty)$.
This closes the final Lean engineering sorry in \texttt{InfiniteZerosBarrier.lean} (the
slit-plane membership gap in the ceml/ceml case), bringing the total machine-verified
theorem count from 24 to \textbf{25}, with exactly 2 genuine open mathematical sorries
remaining --- both requiring o-minimal structure theory.
\end{abstract}

%──────────────────────────────────────────────────────────────────────────────
\section{The Slit-Plane Obstruction}

The proof that every EML tree is real-analytic on $(0,\infty)$ proceeds by lifting
to $\mathbb{C}$: show $x \mapsto t.\mathtt{eval}(\uparrow\!x)$ is $\mathbb{R} \to \mathbb{C}$
analytic, then extract the real part via $\mathtt{reCLM}$.

For a node $\mathtt{ceml}(t_1, t_2)$, the log argument is $t_2.\mathtt{eval}(\uparrow\!x)$.
$\mathtt{Complex.log}$ is $\mathbb{C}$-analytic on the \emph{slit plane}
$\mathbb{C}_{\mathrm{slit}} = \{z \in \mathbb{C} \mid 0 < z.\mathtt{re} \lor z.\mathtt{im} \neq 0\}$
(the plane minus the closed negative real axis, including zero).
For the proof to go through, we need
\[
  t_2.\mathtt{eval}(\uparrow\!x) \in \mathbb{C}_{\mathrm{slit}}
  \quad \forall x \in (0,\infty).
\]
When $t_2$ is a leaf (\texttt{const} or \texttt{var}), this is immediate.
When $t_2$ is itself a $\mathtt{ceml}$ node, it was a sorry.

\subsection{Why the claim is false without extra structure}

\begin{remark}[Counterexample]
Let $t_2 = \mathtt{ceml}(\mathtt{const}\ 0,\ \mathtt{var})$. Then
\[
  t_2.\mathtt{evalReal}(x) = e^0 - \ln x = 1 - \ln x.
\]
For $x > e$, this is negative, hence $t_2.\mathtt{eval}(\uparrow\!x) = 1 - \ln x < 0$
is a non-positive real, \emph{not} in $\mathbb{C}_{\mathrm{slit}}$.
So the slit-plane MapsTo fails without additional conditions on $t_2$.
\end{remark}

%──────────────────────────────────────────────────────────────────────────────
\section{WellFormedPos: The Necessary Condition}

\begin{definition}[WellFormedPos]
An EML tree $t$ is \emph{well-formed at} $x \in (0,\infty)$, written
$t.\mathtt{WellFormedPos}(x)$, if every log argument in $t$ evaluates to a positive real:
\[
\begin{aligned}
  \mathtt{const}\ c,\ x &\mapsto \top \\
  \mathtt{var},\ x &\mapsto \top \\
  \mathtt{ceml}\ t_1\ t_2,\ x &\mapsto t_1.\mathtt{WFP}(x)
    \land t_2.\mathtt{WFP}(x)
    \land 0 < t_2.\mathtt{evalReal}(x).
\end{aligned}
\]
\end{definition}

In Lean 4:
\begin{lstlisting}
def EMLTree.WellFormedPos : EMLTree -> Real -> Prop
  | .const _, _ => True
  | .var,     _ => True
  | .ceml t1 t2, x =>
      t1.WellFormedPos x /\ t2.WellFormedPos x /\ 0 < t2.evalReal x
\end{lstlisting}

\subsection{Why WFP closes the sorry}

Under WFP for the outer node $\mathtt{ceml}\ t_1\ t_2$, we have
$0 < t_2.\mathtt{evalReal}(x)$ as a hypothesis. Since
$t_2.\mathtt{evalReal}(x) = (t_2.\mathtt{eval}(\uparrow\!x)).\mathtt{re}$
by definition, we get $(t_2.\mathtt{eval}(\uparrow\!x)).\mathtt{re} > 0$.
Since $\mathbb{C}_{\mathrm{slit}} = \{z \mid 0 < z.\mathtt{re} \lor z.\mathtt{im} \neq 0\}$,
the left disjunct holds and $t_2.\mathtt{eval}(\uparrow\!x) \in \mathbb{C}_{\mathrm{slit}}$. $\square$

%──────────────────────────────────────────────────────────────────────────────
\section{The Theorems}

\begin{lemma}[eml\_tree\_eval\_analyticOnNhd, 0 sorry]
Let $t : \mathtt{EMLTree}$ and suppose
$\forall x \in (0,\infty),\ t.\mathtt{WellFormedPos}(x)$.
Then $x \mapsto t.\mathtt{eval}(\uparrow\!x)$ is $\mathbb{R} \to \mathbb{C}$ real-analytic
on $(0,\infty)$.

\begin{proof}
By induction on $t$.
\begin{itemize}[noitemsep]
  \item \texttt{const} $c$: constant, analytic trivially.
  \item \texttt{var}: $\mathtt{ofRealCLM}$, a continuous $\mathbb{R}$-linear map, analytic.
  \item $\mathtt{ceml}\ t_1\ t_2$: extract WFP components.
    \begin{itemize}[noitemsep]
      \item \emph{exp part}: $\mathtt{Complex.exp}$ is entire; compose with IH for $t_1$.
      \item \emph{log part}: case split on $t_2$.
        \begin{itemize}[noitemsep]
          \item $\mathtt{const}$: log of a constant is constant, analytic.
          \item $\mathtt{var}$: $\uparrow\!x \in \mathbb{C}_{\mathrm{slit}}$ for $x > 0$
            (by $\mathtt{Complex.ofReal\_mem\_slitPlane}$).
          \item $\mathtt{ceml}$: WFP gives $0 < t_2.\mathtt{evalReal}(x)$,
            hence $(t_2.\mathtt{eval}\ \uparrow\!x).\mathtt{re} > 0$,
            hence $t_2.\mathtt{eval}\ \uparrow\!x \in \mathbb{C}_{\mathrm{slit}}$.
            Apply $\mathtt{hclog.comp}\ (\mathtt{ih2}\ \mathtt{hwf2})\ \mathtt{hmaps}$. \qed
        \end{itemize}
    \end{itemize}
\end{itemize}
\end{proof}
\end{lemma}

\begin{theorem}[\textbf{T\_EML\_WFP\_ANALYTIC}, 0 sorry]\label{thm:main}
Let $t : \mathtt{EMLTree}$ and suppose
$\forall x \in (0,\infty),\ t.\mathtt{WellFormedPos}(x)$.
Then $t.\mathtt{evalReal}$ is real-analytic on $(0,\infty)$.

\begin{proof}
By the lemma above, $x \mapsto t.\mathtt{eval}(\uparrow\!x)$ is analytic.
Since $t.\mathtt{evalReal} = \mathtt{reCLM} \circ (x \mapsto t.\mathtt{eval}(\uparrow\!x))$
(by definition of $\mathtt{evalReal}$), and $\mathtt{reCLM} : \mathbb{C} \to \mathbb{R}$
is a continuous $\mathbb{R}$-linear map (hence analytic), the composition is analytic. $\square$
\end{proof}
\end{theorem}

In Lean 4:
\begin{lstlisting}
lemma eml_tree_analytic (t : EMLTree)
    (hwf : forall x in Set.Ioi 0, t.WellFormedPos x) :
    AnalyticOnNhd R t.evalReal (Set.Ioi 0) := by
  have hcomplex := eml_tree_eval_analyticOnNhd t hwf
  have heq : t.evalReal = Complex.reCLM o (fun x : R => t.eval (x : C)) := by
    ext x; simp [EMLTree.evalReal, Complex.reCLM]
  rw [heq]
  exact (Complex.reCLM.analyticOnNhd Set.univ).comp hcomplex (Set.mapsTo_univ _ _)
\end{lstlisting}

%──────────────────────────────────────────────────────────────────────────────
\section{Relation to Existing Results}

\begin{corollary}[eml\_tree\_analytic\_depth\_le\_1, 0 sorry — no WFP needed]
Every EML tree of depth $\leq 1$ is real-analytic on $(0,\infty)$, without a WFP
hypothesis.
\end{corollary}

\begin{proof}
At depth $\leq 1$, any $\mathtt{ceml}\ t_1\ t_2$ node has $t_2$ of depth 0 (leaf only),
so the ceml/ceml branch is unreachable. Explicit 6-case analysis suffices. $\square$
\end{proof}

The WFP condition is strictly necessary for the general case:
trees with complex constants or nested ceml nodes can produce non-slit-plane
log arguments. T\_EML\_WFP\_ANALYTIC is the correct general statement.

%──────────────────────────────────────────────────────────────────────────────
\section{Remaining Open Sorries}

\textcolor{open}{\textbf{Two genuine open mathematical sorries remain. Both require o-minimal
structure theory (Wilkie 1996).}}

\begin{enumerate}
  \item \textbf{sin\_not\_in\_eml}: $\forall k,\ \forall t : \mathtt{EMLTree},\ t.\mathtt{depth} \leq k \Rightarrow t.\mathtt{evalReal} \neq \sin$. \\
    Requires: every EML-$k$ tree has finitely many zeros on any bounded interval. \\
    This is true because every EML tree is definable in $\mathbb{R}_{\exp}$, which is
    o-minimal (Wilkie, 1996 — not formalized in Mathlib).

  \item \textbf{sin\_not\_in\_real\_EML\_k}: The complex function $z \mapsto \sin(z.\mathtt{re})$
    is not in any $\mathtt{EML\_k}$. Follows from (1) by restriction to $\mathbb{R}$.
\end{enumerate}

Neither sorry is an engineering gap. Both are blocked by the absence of
$\mathbb{R}_{\exp}$ o-minimality in Mathlib.

%──────────────────────────────────────────────────────────────────────────────
\section{Verification Record}

\begin{center}
\begin{tabular}{ll}
\toprule
Item & Value \\
\midrule
Build & \textcolor{proved}{\texttt{lake build → Build completed successfully}} \\
New theorem & T\_EML\_WFP\_ANALYTIC (\texttt{eml\_tree\_analytic}) \\
New predicate & \texttt{EMLTree.WellFormedPos} \\
PROVED\_COUNT & 25 (was 24) \\
Sorry count & 2 (was 3) \\
Closed sorry & ceml/ceml slit-plane MapsTo in \texttt{eml\_tree\_eval\_analyticOnNhd} \\
Remaining sorries & \texttt{sin\_not\_in\_eml}, \texttt{sin\_not\_in\_real\_EML\_k} \\
Commit & e902b2a7 \\
Date & 2026-04-21 \\
\bottomrule
\end{tabular}
\end{center}

\end{document}
